% NOTE: This file is not included by main.tex. The active source is chapters_new.tex and chapters_6_9.tex.
% ==================== 第8章 ====================
\chapter{讨论与局限}
\label{ch:discussion}

\section{主要发现}

\subsection{框架层面}

实验支持三个结论。第一，输出契约应随查询变化。确定问题受益于类型化义务和闭合阅读器；模糊断言受益于锚点、信念和裁决门。第二，模糊分支需要语义抽取器和校准后验两部分。纯LLM锚点后验校准差；纯随机森林后验缺少语义锚点，但在点指标上最强；小权重融合没有提高点指标，因此最终系统保留锚点抽取用于证据资格和审计，点与决策输出使用验证集选中的随机森林后验。第三，路由器在DeepSeek和GLM上以一次短调用恢复几乎全部oracle效用，但8B基座不足以稳定判断语义边界。

\subsection{分支层面}

确定分支的优势来自义务覆盖和置信度回退：它在DeepSeek和GLM上均高于固定16段阅读器，并低于四个具名RAG基线。模糊分支的优势主要来自决策损失和校准，而不是普通准确率；平衡准确率和宏F1的点估计更高，但区间包含零。路由器在较大基座上接近oracle，在8B基座上低于Always-Anchor。

\section{与早期工作的关系}

早期的损失感知策略组合在固定证据模型和固定输出契约上选择固定裁决、弃权或规划器；它是本文模糊分支裁决门的一个特例。本文保留组合作为决策头，但增加了锚点抽取、LLM路由器和确定分支。因而组合是组件，不是完整系统。与把RAG或验证器直接用于所有查询相比，双层框架把“输出什么”与“何时输出”分开处理。

\section{局限与边界}

\subsection{数据与标签}

收入指引实例化是三桶问题；通用十二状态（时机×阶段）只在生产里程碑原型中验证，尚未在大规模基准上评估。标签主要来自SEC XBRL，46条纯文本测试断言被排除在主结果之外；这不是人工金标准，标签噪声是已知边界。FinDVer迁移使用不同阅读器接口，比较是前沿检查而非受控消融。

\subsection{模型与路由}

锚点抽取器会出错，错误硬锚点可能排除正确桶；可靠性门槛和冲突检查只能缓解，不能消除。路由器在构造的混合负载上评估，两个数据集的来源清晰但语义边界可能比真实查询更干净；真实场景中模糊与确定的边界更模糊。Qwen3-8B的DocFinQA精确匹配低于0.11，路由混合效用低于Always-Anchor，说明小基座不适合承担确定阅读或模糊性判断。

% ==================== 第9章 ====================
\chapter{结论}
\label{ch:conclusion}

\section{工作总结}

本文提出FinPlanRAG，一个由LLM路由器连接确定信息分支和模糊断言分支的双层框架。路由器只根据原始查询和截止时点元数据判断模糊性；确定分支绑定类型化义务，按义务覆盖迭代取证，并在证据不足时弃权；模糊分支AnchorEvolve抽取L0/L1/L2锚点，按硬、软、空泛、对抗、过期五类资格判定更新信念，在非对称损失下返回答案、可信集或弃权。两个分支和路由器都调用LLM。

在DocFinQA上，确定分支在DeepSeek上以10.69段取得0.457精确匹配，高于16段固定阅读器的0.390；在GLM上以10.31段取得0.497，高于16段固定阅读器的0.480。在227条XBRL核验收入指引上，模糊分支取得0.533准确率、0.491平衡准确率、0.486宏F1；平衡准确率和宏F1的点估计高于全部六个具名基线，但配对区间包含零；决策损失为3.454，显著低于全部具名基线。路由在DeepSeek和GLM上取得0.995和0.990的准确率，混合效用接近oracle；Qwen3路由较弱，说明基座规模影响语义边界判断。

\section{未来工作}

后续工作有三个方向。第一，把十二状态时机×阶段空间和VoI取证扩展到收入指引之外的生产里程碑、产能投放和监管审批断言。第二，压缩可信集宽度，在保持覆盖率的同时提高信息量。第三，改进路由器对语义边界的判断，例如加入候选段落标题或轻量检索信号，同时保持一次调用的成本。
